\section{环境形式化：POMDP、状态与动作空间}

\subsection{为什么要做形式化}
讲者把 Web Agent 建模为 POMDP，这是非常关键的一步。它让我们可以用强化学习和搜索的统一语言描述问题，而不是把系统行为当作一堆启发式脚本。

\begin{figure}[H]
\centering
\includegraphics[width=0.82\textwidth]{frame-02-pomdp.jpg}
\caption{视频约 00:13:22：讲者将网页交互问题写成 states / actions / observations / transitions}
\end{figure}

\begin{knowledgebox}{POMDP 在这里的直观解释}
\begin{itemize}
    \item \textbf{State}：网页真实状态（包括不可见信息）；
    \item \textbf{Observation}：Agent 当前可见内容（截图、DOM、可访问树等）；
    \item \textbf{Action}：点击、输入、滚动、切换标签、停止等；
    \item \textbf{Reward}：任务是否被正确完成，通常较稀疏。
\end{itemize}
\end{knowledgebox}

\subsection{动作空间设计决定了系统上限}
动作空间过粗会导致表达力不足，过细又会导致搜索空间爆炸。实践中通常需要 ``语义动作 + 低层动作'' 的分层设计，例如 ``搜索餐厅'' 这类子目标由 planner 给出，再由 executor 拆成点击与输入。

\begin{importantbox}{动作空间设计原则}
\begin{enumerate}
    \item 原子动作要可执行且可验证；
    \item 高层动作要能减少搜索深度；
    \item 停止动作必须配套可靠的完成判定；
    \item 每步动作都应带可追踪日志，便于回放和调试。
\end{enumerate}
\end{importantbox}

\begin{table}[H]
\centering
\small
\begin{tabular}{p{3.2cm}p{4.1cm}p{3.7cm}}
\toprule
\textbf{动作层级} & \textbf{典型动作} & \textbf{主要风险} \\
\midrule
低层动作 & click / type / scroll / hover & 步数长、局部错误易累积 \\
中层动作 & open-site / query / filter & 语义歧义、参数绑定困难 \\
高层动作 & complete-subtask / stop & 过早停止或误判完成 \\
\bottomrule
\end{tabular}
\caption{Web Agent 动作层级与风险}
\end{table}

\subsection{奖励稀疏下的训练与评估挑战}
在网页任务里，很多 reward 只在末尾给出，因此信用分配非常困难。系统若只靠终局信号，很难学会中间步骤的精细纠错能力。

\begin{warningbox}{稀疏奖励的副作用}
\begin{itemize}
    \item 训练不稳定，策略更新方差大；
    \item 模型倾向学会投机 shortcut；
    \item 长链路任务的中间错误难以被及时惩罚。
\end{itemize}
\end{warningbox}

\begin{figure}[H]
\centering
\includegraphics[width=0.82\textwidth]{frame-03-trajectory.jpg}
\caption{视频约 00:16:03：讲者展示复杂任务轨迹，强调长链路推理与执行一致性}
\end{figure}

\subsection{本章小结}
POMDP 形式化让 Agent 任务从 ``脚本工程'' 走向 ``决策问题''。但奖励稀疏与动作空间爆炸使得仅靠单次前向推理难以稳定完成复杂任务。

